% Generated draft: requires mathematical and source review.
\documentclass[11pt]{article}
\usepackage[margin=25mm]{geometry}
\usepackage{fontspec}
\setmainfont{Latin Modern Roman}
\usepackage{amsmath,amssymb,mathtools,unicode-math}
\setmathfont{Latin Modern Math}
\usepackage{enumitem,needspace,xurl,hyperref}
\setlength{\parindent}{0pt}
\setlength{\parskip}{5pt}
\setlength{\emergencystretch}{3em}
\newcommand{\N}{\mathbb{N}}
\newcommand{\Z}{\mathbb{Z}}
\newcommand{\Q}{\mathbb{Q}}
\newcommand{\R}{\mathbb{R}}
\newcommand{\C}{\mathbb{C}}
\newcommand{\F}{\mathbb{F}}
\newcommand{\A}{\mathbb{A}}
\newcommand{\PP}{\mathbb{P}}
\newcommand{\E}{\mathbb{E}}
\newcommand{\eps}{\varepsilon}
\newcommand{\dd}{\,\mathrm{d}}
\begin{document}
\section*{UnsolvedMath ID: 30005683}
\textbf{Source problem code:} catalogue entry
\subsection{Definitions and mathematical statement}
Let \(\mu\) denote the diffraction measure of a substitution system and let \(F\) denote its associated Mahler function. The supplied record does not specify the class of substitutions, the weights defining the diffraction measure, its ambient space, or the construction associating \(F\) to the system. Consequently, these symbols describe the intended objects but do not yet determine a quantified mathematical conjecture.

To specify the scaling question, once a Euclidean ambient space and a diffraction measure have been fixed, let \(B_\varepsilon(0)\) be the open ball of radius \(\varepsilon>0\) about the origin and distinguish
\[
M(\varepsilon)=\mu(B_\varepsilon(0)),\qquad
M_*(\varepsilon)=\mu(B_\varepsilon(0)\setminus\{0\}).
\]
These quantities differ when there is an atom at zero. One possible measure of scaling, for either chosen mass function \(H\), is
\[
\alpha_H=\lim_{\varepsilon\downarrow0}
\frac{\log H(\varepsilon)}{\log\varepsilon},
\]
provided the logarithms are defined and the limit exists. This is a possible formalization of ``order of scaling,'' not a convention specified by the source; the source may require a finer asymptotic scale.

For the radial question, the supplied excerpt indicates approach to \(1\) along the real radius, namely \(z=r\to1^-\). A Mahler functional equation relates values of a function at arguments such as \(z,z^q,z^{q^2},\ldots\), for an integer \(q\ge2\); the equation, base, normalization, and relevant asymptotic quantity for \(F\) are not supplied here.

\textbf{Source problem, with unresolved qualifications.} Determine whether, and under which hypotheses, the order of diffraction scaling near zero can be related to the radial asymptotics of the associated Mahler function as \(r\uparrow1\). A precise version must specify the substitution systems under consideration, the correspondence \(\mu\leftrightarrow F\), the treatment of any atom at zero, the meaning of scaling order, and the proposed relation between the two asymptotic quantities. The record supplies no equality, implication, or universal quantifier that can safely be asserted as the intended conjecture.

The supplied dated assessment describes matching exponents in selected Thue--Morse and Rudin--Shapiro cases. The original excerpt also attributes a particular asymptotic to a result of de Bruijn, cited there as reference [4], and says that the speaker's previously mentioned result with Bell does not supply that asymptotic. The displayed formula is corrupted in the supplied transcription, so its precise mathematical content cannot be recovered reliably here.
\subsection{Short English statement}
The question asks whether diffraction mass near zero and the radial behavior of an associated Mahler function encode related asymptotic information. The supplied record labels this question open and mentions selected Thue--Morse and Rudin--Shapiro examples. It does not specify enough hypotheses or a proposed relation to yield a complete conjecture. This is a statement-only draft requiring human mathematical review.
\subsection{Sources}
\begin{itemize}
\item Local UnsolvedMath registry, supplied complete source record: permanent ID \(30005683\), exact source problem code \texttt{catalogue entry}, title \emph{Diffraction Scaling and Mahler-Function Asymptotics}. The record belongs to the Oberwolfach Reports Open Problems collection and gives proposed year 2023; no proposer is identified.
\item Supplied report citation: \emph{Aspects of Aperiodic Order (2024)}, Oberwolfach Reports, pp.\ 2107--2178, DOI \texttt{10.4171/owr/2023/37}. \url{https://doi.org/10.4171/owr/2023/37}. The report title year is reproduced as supplied; the DOI and proposed-year metadata refer to 2023.
\item Supplied literature source labeled ``Baake-Grimm'': \url{https://doi.org/10.1088/1742-5468/2019/05/054003}. No article title, pinpoint reference, or precise supporting result is supplied.
\item The original excerpt attributes an asymptotic to de Bruijn, reference [4], and refers to a result with Bell. The supplied record contains neither the bibliographic entry for [4] nor enough information to identify the Bell result or the speaker.
\item Literature assessment checked at \texttt{2026-08-22}: the supplied curation labels the problem open, describes matching exponents in selected Thue--Morse and Rudin--Shapiro examples, and leaves the general relationship unresolved. This is supplied dated curation, not independently verified literature assessment.
\item The registry labels the statement \texttt{corrected\_verified} and says it was reconstructed and checked against the official report. That is a supplied verification claim; this draft does not independently verify the report, the cited results, or the current problem status.
\end{itemize}
\end{document}
